\section{Local Fourier estimates}\label{sec:local-fourier}

The exterior multiplicity bound in Section~\ref{sec:exterior} uses an
additive density at each prime in $S$.  We now estimate its Fourier transform.
Scalar frequencies detect only the trace and have an exact, order-independent
value.  For the remaining frequencies we record both decay and the support of
the inverse transform on a fixed denominator shell.  The support is needed
because the denominator at a prime dividing the order index can be arbitrarily
large.

\subsection{Quadratic data and the trace marginal}

Fix a rational prime $p$.  We first specify the local hypotheses used in the
argument.  Let
\[
 E=\prod_{\nu=1}^r E_\nu,\qquad R_E=\prod_{\nu=1}^r R_\nu
\]
be a cubic \'etale $\Q_p$-algebra and its maximal order.  Here $E_\nu$ is a
field, $R_\nu$ is its integer ring, and $r\le3$.  Write
$\mathfrak q_\nu=(\pi_\nu)$ for the maximal ideal of $R_\nu$,
$v_\nu(\pi_\nu)=1$, $e_\nu=v_\nu(p)$, and
$q_\nu=|R_\nu/\mathfrak q_\nu|=p^{f_\nu}$.  Thus
\[
 \sum_\nu e_\nu f_\nu=3,
 \qquad |\Nm_{E_\nu/\Q_p}a|_p=q_\nu^{-v_\nu(a)}.
\]
Let $\mathfrak d_\nu=\mathfrak q_\nu^{d_\nu}$ be the different of
$E_\nu/\Q_p$, and put $\mathfrak d_E=\prod_\nu\mathfrak d_\nu$.

Let $W$ be a free $E$-module of rank two, with its unital action denoted by
$\rho:E\longrightarrow\End_{\Q_p}(W)$.  Suppose that $B$ is a nondegenerate
symmetric $\Q_p$-bilinear form satisfying
\[
 B(\rho(a)x,y)=B(x,\rho(a)y)\qquad(a\in E,\ x,y\in W).
\]
Let $L$ be a full $\Z_p$-lattice in $W$ which is self-dual for $B$:
\[
 L=\{y\in W:B(y,L)\subset\Z_p\}.
\]
Assume also that $Q_1(x)=B(x,x)/2$ is integral on $L$.  For the scalar
calculation we shall use the stronger property that in some integral
coordinates on $L$,
\begin{equation}\label{eq:fourier-split-form}
 Q_1(x_1,y_1,x_2,y_2,x_3,y_3)
       =x_1y_1-x_2y_2+x_3y_3.
\end{equation}
The division by $2$ takes place in $\Q_p$.  Equation~\eqref{eq:fourier-split-form}
asserts the required integrality also when $p=2$.  The shell estimates below
use only $Q_1(L)\subset\Z_p$; the exact split expression is used for the
scalar marginal.

For $a\in E$, set $Q_a(x)=B(\rho(a)x,x)/2$.  The nondegenerate trace pairing
on $E$ determines a unique quadratic map $\Psi_E:W\longrightarrow E$ by
\begin{equation}\label{eq:fourier-polar-identity}
 \Tr_{E/\Q_p}(a\Psi_E(x))=Q_a(x)
       =\frac12 B(\rho(a)x,x)\qquad(a\in E).
\end{equation}
Define the stabilizer order and the auxiliary order
\[
 \mathcal C_L=\{a\in E:\rho(a)L\subset L\},
 \qquad \mathcal C_L^\sharp=\Z_p+2\mathcal C_L.
\]
Indeed $\mathcal C_L$ is an order in $R_E$.  It is a subring containing $1$,
and a sufficiently large power of $p$ times $R_E$ preserves $L$.  Conversely,
an element preserving $L$ satisfies the monic characteristic polynomial of
its integral matrix on $L$.  The action is faithful, so the element is
integral over $\Z_p$ and belongs to $R_E$.

In the application, $E=F_p$, $W=V_p$, $L=L_{V,p}$, $B=B_V$, and
$\mathcal C_L=C_{F,p}$.  Lemma~\ref{lem:ext-quadratic-structure} supplies
all the stated hypotheses and identifies $\Psi_E$ with the base change of
$\Psi$.  We use the abstract notation temporarily to include every product
algebra and every ramification type in one argument.

Choose an additive character $\psi_p$ of $\Q_p$ with conductor $\Z_p$.
We give $W$ the self-dual Haar measure for $(x,y)\mapsto\psi_p(B(x,y))$;
it gives $L$ mass $1$.  On $E$ use the self-dual measure for
$(A,z)\mapsto\psi_p(\Tr_{E/\Q_p}(Az))$, denoted by $dz$ or $dA$ according
to the variable.  If $\mathfrak a$ is a fractional $R_E$-ideal, its
annihilator for this pairing is
$\mathfrak a^\perp=\mathfrak a^{-1}\mathfrak d_E^{-1}$.  In particular
\begin{equation}\label{eq:fourier-self-dual-volume}
 \vol(R_E)=\prod_\nu q_\nu^{-d_\nu/2}\le1.
\end{equation}
To see the normalization, $R_E^\perp=\mathfrak d_E^{-1}$ and
$\vol(R_E)\vol(R_E^\perp)=1$, whereas
$[\mathfrak d_E^{-1}:R_E]=\prod_\nu q_\nu^{d_\nu}$.
For $E=F_p$, Equation~\eqref{eq:fourier-self-dual-volume} reads
$\vol(R_{F,p})=|D_F|_p^{1/2}$.

Let $\mu_p$ be the law of
\begin{equation}\label{eq:fourier-density-law}
 u\Psi_E(x),\qquad x\in L\text{ with additive probability},
 \qquad u\in\Z_p^\times\text{ with Haar probability},
\end{equation}
with $x$ and $u$ independent.  This is the law defining $f_p^*$ in
Lemma~\ref{lem:ext-density}.  The following regularity fact specifies the
version of that density which we shall invert pointwise.

\begin{lemma}\label{lem:fourier-density}
The measure $\mu_p$ has a density $f_p^*$ with respect to $dz$.  It has a
finite locally constant version at every $z=(z_\nu)\in E$ with
$z_\nu\ne0$ for all $\nu$.  This version vanishes outside
\[
 (\mathcal C_L^\sharp)^\vee
 =\{z\in E:\Tr_{E/\Q_p}(z\mathcal C_L^\sharp)\subset\Z_p\}.
\]
In particular $\Tr_{E/\Q_p}(u\Psi_E(x))\in\Z_p$ on the domain of
Equation~\eqref{eq:fourier-density-law}.
\end{lemma}

\begin{proof}
Define an $E$-valued bilinear form $b_E$ by
\[
 \Tr_{E/\Q_p}(a b_E(x,y))=B(\rho(a)x,y)\qquad(a\in E).
\]
Self-adjointness of the action makes $b_E$ symmetric and $E$-bilinear.
It is nondegenerate because $B$ is nondegenerate.  Equation~\eqref{eq:fourier-polar-identity}
then gives $\Psi_E(x)=b_E(x,x)/2$ and
$d\Psi_{E,x}(h)=b_E(x,h)$.  On the $\nu$-th field factor, a nonzero
component of $\Psi_E(x)$ implies that the component $x_\nu$ of $x$ is
nonzero.  The nondegenerate form on the rank-two $E_\nu$-space makes
$h_\nu\mapsto b_E(x_\nu,h_\nu)$ a nonzero, hence onto, $E_\nu$-linear
functional.  It follows that $\Psi_E$ is a $\Q_p$-submersion at every
preimage of an output with nonzero components.  Multiplication by a unit
$u$ is invertible on $E$, so the map $(x,u)\mapsto u\Psi_E(x)$ is a
submersion at the same points.

For each $\nu$ the polynomial $\Psi_{E,\nu}$ is nonzero: a nondegenerate
symmetric form in characteristic zero has a nonzero quadratic form.  Choose
a nonzero $\Q_p$-coordinate polynomial of this component.  Its zero set in
the open lattice $L$ has Haar measure zero, by the usual one-variable root
bound followed by induction and Fubini.  Off the union of
these null sets, the local submersion theorem gives absolute continuity of
the pushforward.  For a fixed output with nonzero components, its preimage
in the compact set $L\times\Z_p^\times$ is compact.  Cover this preimage
by finitely many compact-open submersion charts, refined to disjoint
compact-open pieces within those charts.  The image of the remaining compact
set avoids a sufficiently small neighborhood of the output.  After shrinking
that neighborhood, integration in the fiber coordinates of the pieces gives
a finite constant density on it.  This proves the asserted
pointwise version, including value zero when the neighborhood misses the
image.

Finally, for $c\in\mathcal C_L$ and $r\in\Z_p$,
\[
 \Tr_{E/\Q_p}((r+2c)\Psi_E(x))
       =rQ_1(x)+B(\rho(c)x,x)\in\Z_p\qquad(x\in L).
\]
Multiplication by $u\in\Z_p^\times$ preserves this inclusion.  The image
of the compact domain is therefore contained in the displayed trace dual,
which proves the support and trace assertions.
\end{proof}

We use the Fourier convention
\[
 \widehat f(A)=\int_E f(z)\psi_p(\Tr_{E/\Q_p}(Az))\,dz,
 \qquad
 \check g(z)=\int_E g(A)\psi_p(-\Tr_{E/\Q_p}(Az))\,dA.
\]
The transform of the measure in Equation~\eqref{eq:fourier-density-law} is
\begin{equation}\label{eq:fourier-gamma}
 \gamma_p(A)=\widehat{f_p^*}(A)
 =\E_{u\in\Z_p^\times}\int_L\psi_p(uQ_A(x))\,dx.
\end{equation}

\begin{lemma}[The scalar trace density]\label{lem:fourier-scalar-density}
Suppose that the integral scalar form is exactly
$Q_1=x_1y_1-x_2y_2+x_3y_3$ as in
Equation~\eqref{eq:fourier-split-form}.  Regard $a\in\Q_p$ as a scalar in
$E$.  In the quartic application this scalar form is $Q_1=Q_V$.  Then, for
every prime $p$, including $p=2$,
\begin{equation}\label{eq:fourier-scalar-coefficients}
 \gamma_p(a)=
 \begin{cases}
 1,&a\in\Z_p,\\
 p^{-3k},&v_p(a)=-k<0.
 \end{cases}
\end{equation}
The trace of the law in Equation~\eqref{eq:fourier-density-law} has density
$d_p(t)\mathbf1_{\Z_p}(t)$ for additive measure giving $\Z_p$ mass $1$.
This density is continuous on $\Z_p$ and satisfies
\begin{equation}\label{eq:fourier-scalar-density-bound}
 |d_p(t)-1|\le\frac1{p(p+1)}\qquad(t\in\Z_p).
\end{equation}
In particular $5/6\le d_p(t)\le7/6$, and for every finite set of primes
$S$ and every choice $t_p\in\Z_p$,
\[
 \prod_{p\in S}d_p(t_p)\asymp1
\]
with absolute upper and lower constants.
\end{lemma}

\begin{proof}
For fixed $u\in\Z_p^\times$ and $v_p(a)=-k<0$, integration in the second
coordinate gives
\[
 \int_{\Z_p^2}\psi_p(ua xy)\,dx\,dy
 =\int_{\Z_p}\mathbf1_{\{ua x\in\Z_p\}}\,dx=p^{-k}.
\]
For $a\in\Z_p$ the integral is $1$.  The sign of $a$ has no effect, so the
three independent pairs in Equation~\eqref{eq:fourier-split-form} give
Equation~\eqref{eq:fourier-scalar-coefficients}, even before averaging $u$.
This calculation uses the integral split expression and does not divide by
$2$ in $\Z_p$.

The scalar coefficient is the Fourier coefficient of the trace law, since
$\Tr(a z)=a\Tr(z)$.  That law is supported on $\Z_p$ by
Lemma~\ref{lem:fourier-density}.  The character group of $\Z_p$ is
$\Q_p/\Z_p$, and it has $p^k-p^{k-1}$ characters of exact denominator
$p^k$.  Hence its nonconstant Fourier coefficients have total absolute sum
\[
 \sum_{k\ge1}(p^k-p^{k-1})p^{-3k}
       =\frac{1-p^{-1}}{p^2-1}=\frac1{p(p+1)}.
\]
Fourier inversion on the compact group $\Z_p$ is therefore uniformly
convergent and gives the continuous density and
Equation~\eqref{eq:fourier-scalar-density-bound}.  The lower bound is
positive at every prime, and $\sum_p1/(p(p+1))<\infty$.  Taking logarithms
of the upper and lower bounds proves the last assertion.
\end{proof}

\subsection{Denominator shells and quadratic transformation}

For $A=(A_\nu)\in E$, its denominator ideal is
\[
 \mathfrak h_p(A)=\prod_\nu\mathfrak q_\nu^{h_\nu(A)},\qquad
 h_\nu(A)=\max\{0,-v_\nu(A_\nu)\},
\]
where $v_\nu(0)=+\infty$.  For an integral $R_E$-ideal
$\mathfrak h_p=\prod_\nu\mathfrak q_\nu^{h_\nu}$, its shell is thus
\[
 \mathscr S_{\mathfrak h_p}
 =\prod_{h_\nu=0}R_\nu\ \times\!
  \prod_{h_\nu>0}\pi_\nu^{-h_\nu}R_\nu^\times.
\]
The integral components of a shell may be zero.  Define
\begin{equation}\label{eq:fourier-shell-indices}
 H_p=\Nm(\mathfrak h_p)=\prod_\nu q_\nu^{h_\nu},\qquad
 p^{g_p}=[\mathfrak h_p:\mathfrak h_p\cap\mathcal C_L].
\end{equation}
The shell is a compact open subset of $\mathfrak h_p^{-1}$.  We shall shift
its frequencies by integral scalars so that the quadratic form is
nondegenerate, without changing the integral in
Equation~\eqref{eq:fourier-gamma}.

\begin{lemma}[A scalar shift on a shell]\label{lem:fourier-shift}
Put $\iota_p=1$ for $p\le5$ and $\iota_p=0$ for $p>5$.  For every
$A\in\mathscr S_{\mathfrak h_p}$ there is an $s\in\{0,1,2,3\}$ such that
$\alpha=A+s\in E^\times$ and
\[
 v_\nu(\alpha_\nu)=-h_\nu\quad(h_\nu>0),\qquad
 0\le v_\nu(\alpha_\nu)\le e_\nu\iota_p\quad(h_\nu=0).
\]
Writing $\beta=\alpha^{-1}$, one has
\begin{equation}\label{eq:fourier-shift-bounds}
 \beta\in p^{-\iota_p}\mathfrak h_p,
 \qquad p^{-3\iota_p}H_p
       \le |\Nm_{E/\Q_p}\alpha|_p\le H_p.
\end{equation}
In particular the norm is exactly $H_p$ when $p>5$.  If
\begin{equation}\label{eq:fourier-j-definition}
 j_p=j_p(A,s)=\min\{j\ge0:p^j\beta\in\mathcal C_L\},
\end{equation}
then $j_p$ is finite and $j_p\le g_p+\iota_p$.  Finally, for any generator
$i\in E^\times$ of the ideal $\mathfrak h_p$, the element $x=i(A+s)$ is in
$R_E$, is a unit of $R_E$ if $p>5$, and satisfies
$x^{-1}\in p^{-\iota_p}R_E$.
\end{lemma}

\begin{proof}
An integral scalar shift leaves every pole valuation $-h_\nu<0$ unchanged.
For an integral component and $p>5$, the condition that $A_\nu+s$ fail to be
a unit forbids at most one residue $s\in\mathbf F_p$: it forbids none if
the residue of $-A_\nu$ does not belong to the scalar subfield.  There are
at most three components, and $0,1,2,3$ are distinct modulo $p$.  Some choice
therefore makes all integral components units.

For $p\le5$, normalize the extension of $v_p$ by $v_p(p)=1$, so that on
$E_\nu$ it is $v_\nu/e_\nu$.  Distinct elements $s,s'$ of the list satisfy
$v_p(s-s')\le1$.  The ultrametric inequality shows that for a fixed integral
component at most one candidate can satisfy $v_p(A_\nu+s)>1$.  Again there
are at most three components, so a candidate remains for which all their
valuations are at most $1$.  Thus $A_\nu+s$ is nonzero in every integral
component.

The bounds on $\beta$ follow componentwise.  In the norm, the possible
loss from the integral components is at most
$p^{-\iota_p\sum_\nu e_\nu f_\nu}=p^{-3\iota_p}$, proving
Equation~\eqref{eq:fourier-shift-bounds}.  The group
$\mathfrak h_p/(\mathfrak h_p\cap\mathcal C_L)$ has order $p^{g_p}$ and is
killed by $p^{g_p}$.  Thus
\begin{equation}\label{eq:fourier-order-containment}
 p^{g_p}\mathfrak h_p\subset\mathcal C_L.
\end{equation}
Multiplying the containment for $\beta$ by $p^{g_p+\iota_p}$ proves the
claim about $j_p$.  Finally $v_\nu(i_\nu)=h_\nu$.  The components of
$i\alpha$ consequently have valuation zero on the pole factors and valuation
between zero and $e_\nu\iota_p$ on the others.  This proves the assertions
about $x$ and $x^{-1}$.
\end{proof}

For the shift in Lemma~\ref{lem:fourier-shift},
$Q_{A+s}=Q_A+sQ_1$ and $Q_1(L)\subset\Z_p$.  Thus for every fixed unit $u$
\begin{equation}\label{eq:fourier-scalar-shift-invariance}
 \int_L\psi_p(uQ_{A+s}(x))\,dx
       =\int_L\psi_p(uQ_A(x))\,dx.
\end{equation}
The following quadratic transformation turns the norm of $A+s$ into decay,
and leaves an integral involving its inverse.

\begin{lemma}[Quadratic Fourier transformation]\label{lem:fourier-quadratic-transform}
Let $W_0$ be a finite-dimensional $\Q_p$-space with a nondegenerate symmetric
form $B_0$, and let $L_0\subset W_0$ be a full $\Z_p$-lattice self-dual for
$B_0$.  Use the self-dual Haar
measure, which gives $L_0$ mass $1$.  For an invertible $B_0$-self-adjoint
operator $P$, put $q_P(x)=B_0(Px,x)/2$.  There is a phase $\omega(P)$ with
$|\omega(P)|=1$ such that
\begin{equation}\label{eq:fourier-general-quadratic}
 \int_{L_0}\psi_p(q_P(x))\,dx
 =\omega(P)|\det P|_p^{-1/2}
       \int_{L_0}\psi_p(-q_{P^{-1}}(x))\,dx.
\end{equation}
The normalized phase is invariant under linear isometry of the quadratic
form: $\omega(T^*PT)=\omega(P)$ for every invertible $T$, where $T^*$ is
the $B_0$-adjoint.

For the local data above, $\alpha\in E^\times$, $\beta=\alpha^{-1}$ and
$u\in\Z_p^\times$, this gives
\begin{equation}\label{eq:fourier-quadratic-transform}
 \int_L\psi_p(uQ_\alpha(x))\,dx
 =\omega_p(u,\alpha)|\Nm_{E/\Q_p}\alpha|_p^{-1}
       \int_L\psi_p(-u^{-1}Q_\beta(x))\,dx,
 \qquad |\omega_p(u,\alpha)|=1.
\end{equation}
For every $r\in E^\times$ one has
$\omega_p(u,\alpha r^2)=\omega_p(u,\alpha)$.
All assertions hold at $p=2$.
\end{lemma}

\begin{proof}
We use Weil's Fourier formula for a quadratic character
\cite[Chapter~I, \S14, Theorem~2, p.~161]{Weil1964}.
We record its normalization for the transform
\[
 \mathcal F_{B_0}h(y)=\int h(x)\psi_p(B_0(x,y))\,dx.
\]
Completing the square gives
\[
 q_P(x)+B_0(x,y)=q_P(x+P^{-1}y)-q_{P^{-1}}(y).
\]
Let $B_N=p^{-N}L_0$ and $c_N=\int_{B_N}\psi_p(q_P(x))\,dx$.  On any fixed
compact set of $y$, translation by $P^{-1}y$ preserves $B_N$ for all
sufficiently large $N$.  The transform of
$\mathbf1_{B_N}\psi_p(q_P)$ on that set is therefore
$c_N\psi_p(-q_{P^{-1}})$.  If this identity is paired with a
Schwartz--Bruhat function, its left side stabilizes with $N$, since the
Fourier transform of the test function has compact support.  Choosing a
test function whose integral against $\psi_p(-q_{P^{-1}})$ is nonzero shows
that $c_N$ stabilizes as well.  This proves the distributional shape of the
formula.

Differencing the absolute square of $c_N$ and integrating in one variable
gives
\[
 |c_N|^2=\vol(B_N)
   \int_{B_N\cap P^{-1}B_N^\perp}\psi_p(q_P(h))\,dh.
\]
For sufficiently large $N$, the intersection is $P^{-1}p^NL_0$ and
$q_P$ is integral on it.  The right side then equals
$\vol(p^{-N}L_0)\vol(P^{-1}p^NL_0)=|\det P|_p^{-1}$.
Hence the stabilized constant is $\omega(P)|\det P|_p^{-1/2}$ with unit
phase.  The transform of $\mathbf1_{L_0}$ is $\mathbf1_{L_0}$, so pairing
the distribution identity with this function gives
Equation~\eqref{eq:fourier-general-quadratic}.

Under $q_{T^*PT}=q_P\circ T$, changing variables multiplies the Gaussian
constant by $|\det T|_p^{-1}$ and its normalizing determinant square root by
$|\det T|_p$.  Their product, the phase, is unchanged; this is also the
isometry invariance in Weil~\cite[Chapter~II, \S25, p.~173]{Weil1964}.
For $P=u\rho(\alpha)$ its inverse is $u^{-1}\rho(\beta)$.  Rank two over
$E$ gives $\det\rho(\alpha)=\Nm_{E/\Q_p}(\alpha)^2$, and $|u|_p=1$.
This proves Equation~\eqref{eq:fourier-quadratic-transform}.  Taking
$T=\rho(r)$ proves the last phase identity.  The polar operator of $q_P$ is
$P$, because its polarization is $B_0(Px,y)$.  In particular this proof
never assumes that $2$ is a unit of $\Z_p$.
\end{proof}

\subsection{The inverse transform on a shell}

Extend the restriction of $\gamma_p$ to a shell by zero, and denote its
inverse transform by
\begin{equation}\label{eq:fourier-shell-transform}
 \gamma_{p,\mathfrak h_p}(A)
   =\mathbf1_{\mathscr S_{\mathfrak h_p}}(A)\gamma_p(A),
 \qquad
 K_{p,\mathfrak h_p}(z)
   =\int_E\gamma_{p,\mathfrak h_p}(A)
             \psi_p(-\Tr_{E/\Q_p}(Az))\,dA.
\end{equation}
The next lemma gives the estimates for these functions which will be used
in Poisson summation.  Its quantitative support assertions separate the case where
the denominator ideal is contained in the stabilizer order.

\begin{lemma}[Intrinsic shell estimates]\label{lem:fourier-shell-estimates}
Let $E$ be the cubic algebra above, let $W$ have rank two over $E$ with
self-adjoint action for a nondegenerate symmetric form $B$, and let $L$ be
self-dual for $B$, with $Q_1(L)\subset\Z_p$.  Use the stabilizer order
$\mathcal C_L=\{a\in E:\rho(a)L\subset L\}$ and $\Psi_E,\gamma_p$ from
Equations~\eqref{eq:fourier-polar-identity} and \eqref{eq:fourier-gamma}.
Fix an integral ideal
$\mathfrak h_p=\prod_\nu\mathfrak q_\nu^{h_\nu}$, and use
$H_p,g_p$ from Equation~\eqref{eq:fourier-shell-indices}.
The zero-extended function $\gamma_{p,\mathfrak h_p}$ and its inverse
$K_{p,\mathfrak h_p}$ are Schwartz--Bruhat functions.  Moreover,
\[
 \supp K_{p,\mathfrak h_p}\subset(\mathcal C_L^\sharp)^\vee
 \subset\{z\in E:\Tr_{E/\Q_p}z\in\Z_p\}.
\]
For
$A\in\mathscr S_{\mathfrak h_p}$ choose $s$ and $j_p=j_p(A,s)$ as in
Lemma~\ref{lem:fourier-shift}.  Then
\begin{equation}\label{eq:fourier-shell-decay}
 |\gamma_p(A)|\le p^{3\iota_p}H_p^{-1}p^{-j_p/2},
 \qquad
 \|K_{p,\mathfrak h_p}\|_\infty
       \le p^{3\iota_p}\vol(R_E)\le p^{3\iota_p}.
\end{equation}
Here $\iota_p=\mathbf1_{\{p\le5\}}$, so the first constant is $1$ at
every $p>5$ and is at most $125$ at the three remaining primes.

If $g_p>0$ or $p\le5$, then
\begin{equation}\label{eq:fourier-inverse-support}
 \supp K_{p,\mathfrak h_p}
       \subset p^{-4(g_p+1)}\mathfrak h_p\mathfrak d_E^{-1}.
\end{equation}
If $g_p=0$ and $p>5$, then support instead requires, for each $\nu$,
\begin{equation}\label{eq:fourier-boundary-support}
 z_\nu\in
 \begin{cases}
 \mathfrak d_\nu^{-1},&h_\nu=0,\\
 \mathfrak q_\nu^{h_\nu-1}\mathfrak d_\nu^{-1},&h_\nu>0.
 \end{cases}
\end{equation}
For a $z$ satisfying these support conditions, define its set of boundary
primes by
\[
 \partial_{\mathfrak h_p}(z)
 =\{\mathfrak q_\nu:h_\nu>0,
       \ z_\nu\notin\mathfrak q_\nu^{h_\nu}\mathfrak d_\nu^{-1}\}.
\]
In this case
\begin{equation}\label{eq:fourier-boundary-decay}
 |K_{p,\mathfrak h_p}(z)|
 \le 8\vol(R_E)
       \prod_{\mathfrak q_\nu\in\partial_{\mathfrak h_p}(z)}q_\nu^{-1/2}.
\end{equation}
Every constant is uniform in the lattice, the order $\mathcal C_L$, the
shell, and the field factors and their ramification.
\end{lemma}

\begin{proof}
For $c\in\mathcal C_L^\sharp$, Equation~\eqref{eq:fourier-polar-identity}
shows that $Q_c(L)\subset\Z_p$.  Consequently $\gamma_p(A+c)=\gamma_p(A)$.
Each denominator shell is invariant under translation by $R_E$: adding an
integral element preserves every negative component valuation and preserves
the set of integral components.  Since $\mathcal C_L^\sharp\subset R_E$,
the zero-extended shell restriction retains this period.  Its inverse is
therefore supported in the annihilator $(\mathcal C_L^\sharp)^\vee$, whose
elements have integral trace because $1\in\mathcal C_L^\sharp$.
The order $\mathcal C_L^\sharp$ is an open lattice, and the shell is compact
open.  Thus the zero extension in Equation~\eqref{eq:fourier-shell-transform}
is Schwartz--Bruhat, as is its Fourier inverse.

For a fixed unit $u$ put
$I_\beta(u)=\int_L\psi_p(-u^{-1}Q_\beta(x))\,dx$.
Polarization and additive orthogonality on the self-dual lattice give
\[
 |I_\beta(u)|^2
 =\int_{\substack{h\in L\\\rho(\beta)h\in L}}
       \psi_p(-u^{-1}Q_\beta(h))\,dh.
\]
Indeed, in the difference $Q_\beta(y+h)-Q_\beta(y)$ the term linear in $y$
is $B(\rho(\beta)h,y)$, whose integral is the indicator of
$\rho(\beta)h\in L$.  The absolute value of the last display is at most the
relative mass of $L\cap\rho(\beta)^{-1}L$.  If
$p^{a_1},\ldots,p^{a_6}$ are the Smith elementary divisors of
$\rho(\beta)$ on $L$, then
\[
 j_p=\max\{0,-\min_i a_i\},\qquad
 [L:L\cap\rho(\beta)^{-1}L]
       =p^{\sum_i\max(0,-a_i)}\ge p^{j_p}.
\]
The equality for $j_p$ is exactly the definition of $\mathcal C_L$.
It follows that $|I_\beta(u)|\le p^{-j_p/2}$.  Apply
Equations~\eqref{eq:fourier-scalar-shift-invariance} and
\eqref{eq:fourier-quadratic-transform}, then use
Equation~\eqref{eq:fourier-shift-bounds} and average over $u$.  This proves
the first bound in Equation~\eqref{eq:fourier-shell-decay}.  The shell is
contained in $\mathfrak h_p^{-1}$, which has volume
$H_p\vol(R_E)$.  Integration of the first bound, with $p^{-j_p/2}\le1$,
proves the second.

We next prove Equation~\eqref{eq:fourier-inverse-support}.  Put
$N=4(g_p+1)$ and take $t\in p^N\mathfrak h_p^{-1}$.  For $A$ in the shell
choose $s$ as in Lemma~\ref{lem:fourier-shift}, and write
$\alpha=A+s$, $\beta=\alpha^{-1}$.  The component valuations of $t$ are
at least $-h_\nu+e_\nu N$.  Translation by $t$ therefore preserves the
shell, and the same $s$ still has the valuation properties of that lemma.
Set $\zeta=t\beta$.  Equation~\eqref{eq:fourier-shift-bounds} gives
\[
 \zeta\in p^{N-\iota_p}R_E,\qquad
 \alpha+t=\alpha(1+\zeta),\qquad
 (\alpha+t)^{-1}-\beta
       =-\frac{t\beta^2}{1+\zeta}
       \in p^{N-2\iota_p}\mathfrak h_p.
\]
Here $1+\zeta$ is a unit.  It is a square unit in each factor when $p$ is
odd, by Hensel's lemma on principal units.  When $p=2$, one has
$\zeta\in8R_E$.  In the $\nu$-th field the Hensel inequality for
$Y^2-(1+\zeta_\nu)$ at $Y=1$ is
$v_\nu(\zeta_\nu)\ge3e_\nu>2e_\nu=2v_\nu(2)$, so it is a square unit
there as well, including for ramified dyadic factors.

Moreover $N-2\iota_p\ge g_p+v_p(2)$.  Equation~\eqref{eq:fourier-order-containment}
therefore places the inverse difference in $2\mathcal C_L$.  Its $Q$-value
on $L$ is integral by Equation~\eqref{eq:fourier-polar-identity}.
In Equation~\eqref{eq:fourier-quadratic-transform}, the inverse integral is
unchanged by this difference.  The norm is unchanged by the unit $1+\zeta$,
and the phase is unchanged because that unit is a square.  The original
integral, for each fixed $u$, is consequently unchanged under $A\mapsto A+t$.
So is its unit average.  Translation by $p^N\mathfrak h_p^{-1}$ preserves
the shell as a set, hence preserves its zero-extended function.  The inverse
transform is supported on the annihilator of this group, namely
$p^{-N}\mathfrak h_p\mathfrak d_E^{-1}$.  This proves
Equation~\eqref{eq:fourier-inverse-support} with the absolute exponent $4$.

It remains to prove the sharper assertions for $g_p=0$ and $p>5$.  The
equality $g_p=0$ says $\mathfrak h_p\subset\mathcal C_L$.  For each factor
with $h_\nu=0$, its component idempotent $1_\nu$ and its whole component
ring $1_\nu R_\nu$ lie in $\mathfrak h_p$.  The self-adjoint projections
$\rho(1_\nu)$ preserve $L$ and split off their lattices orthogonally.
The complementary projection $1-\sum_{h_\nu=0}\rho(1_\nu)$ also preserves
$L$.  Each resulting summand is self-dual, because the splitting is
orthogonal and $L$ is self-dual.

On an $h_\nu=0$ summand, an integral $A_\nu$ preserves its lattice, so
$Q_{A_\nu}$ is integral there since $2$ is a unit.  That factor of the
integral in Equation~\eqref{eq:fourier-gamma} is $1$.  On the combined
summand of all $h_\nu>0$, the frequency $A_+$ is invertible.  Its inverse,
extended by zero on the other factors, lies in
$\mathfrak h_p\subset\mathcal C_L$.  Its inverse quadratic integral in
Equation~\eqref{eq:fourier-general-quadratic} is therefore also $1$.
On that summand the square root of the polar determinant has absolute value
$H_p$, because each positive factor still has rank two.  Applying the formula
therefore shows that the remaining frequency function is $H_p^{-1}$ times a
unit-average of Gaussian phases.  Its absolute value is at most $H_p^{-1}$.
If there are no positive
$h_\nu$, this conclusion has $H_p=1$ and the empty summand contributes $1$.

Let $\chi_\nu$ be the quadratic character of
$(R_\nu/\mathfrak q_\nu)^\times$ for a factor with $h_\nu>0$.  At odd
residue characteristic every principal unit is a square, so unit square
classes are exactly the two residue square classes.  Phase invariance in
Lemma~\ref{lem:fourier-quadratic-transform} shows that the bounded unit-average
just obtained depends only on the tuple of these classes.  This use of phase
invariance concerns the ambient quadratic form; a unit need not preserve
the lattice.  Fourier expansion on the product of the two-element
square-class groups expresses the shell function as a sum of at most
$2^{|\{\nu:h_\nu>0\}|}\le8$ terms of the form
\[
 c_{\boldsymbol\varepsilon}H_p^{-1}
 \prod_{h_\nu=0}\mathbf1_{R_\nu}(A_\nu)
 \prod_{h_\nu>0}
   \mathbf1_{\pi_\nu^{-h_\nu}R_\nu^\times}(A_\nu)
   \chi_\nu(\overline{\pi_\nu^{h_\nu}A_\nu})^{\varepsilon_\nu},
 \qquad |c_{\boldsymbol\varepsilon}|\le1,
\]
where $\varepsilon_\nu\in\{0,1\}$.  The bound for the coefficients follows
from the normalized finite Fourier expansion of a function of absolute
value at most $1$.  It also covers any coupling of components caused by
the common average over $u$.

We compute the inverse of each factor.  Put
$\psi_\nu=\psi_p\circ\Tr_{E_\nu/\Q_p}$ and use the component self-dual
measure.  For $h_\nu=0$, the integral over $R_\nu$ is
$\vol(R_\nu)\mathbf1_{\mathfrak d_\nu^{-1}}(z_\nu)$.
For $h=h_\nu>0$, the normalized inverse of the trivial unit character is
\[
 q_\nu^{-h}\int_{\pi_\nu^{-h}R_\nu^\times}
       \psi_\nu(-Az)\,dA
 =\vol(R_\nu)\left(
    \mathbf1_{\mathfrak q_\nu^h\mathfrak d_\nu^{-1}}(z)
   -q_\nu^{-1}\mathbf1_{\mathfrak q_\nu^{h-1}\mathfrak d_\nu^{-1}}(z)
                     \right).
\]
This follows by subtracting the integral over $\pi_\nu^{-h+1}R_\nu$ from
that over $\pi_\nu^{-h}R_\nu$ and using their annihilators.  For the quadratic
unit character, its function of $a=\pi_\nu^hA$ is invariant under addition
by $\mathfrak q_\nu$.  This forces the normalized inverse to vanish outside
$\mathfrak q_\nu^{h-1}\mathfrak d_\nu^{-1}$.  It also vanishes on
$\mathfrak q_\nu^h\mathfrak d_\nu^{-1}$, where the additive character on
$R_\nu$ is trivial and the sum of $\chi_\nu$ is zero.  At the intervening
boundary the additive character descends to a nontrivial character $\eta$
of the residue field $k_\nu$, and the normalized integral is
$\vol(R_\nu)q_\nu^{-1}$ times its quadratic Gauss sum.

For completeness, if $\chi$ is the quadratic character of a finite field
$k$ of odd order $q$ and $\eta$ is a nontrivial additive character, then
\[
 \left|\sum_{a\in k^\times}\chi(a)\eta(a)\right|^2
 =\sum_{t\in k^\times}\chi(t)
       \sum_{b\in k^\times}\eta((t-1)b)
 =(q-1)-\sum_{\substack{t\in k^\times\\t\ne1}}\chi(t)=q.
\]
Thus the quadratic factor has magnitude
$\vol(R_\nu)q_\nu^{-1/2}$ at the boundary.  The trivial factor there has
the smaller magnitude $\vol(R_\nu)q_\nu^{-1}$.  Multiplying these
component estimates, using $H_p=\prod_{h_\nu>0}q_\nu^{h_\nu}$ and
$\prod_\nu\vol(R_\nu)=\vol(R_E)$, and summing at most eight terms proves
Equations~\eqref{eq:fourier-boundary-support} and
\eqref{eq:fourier-boundary-decay}.  The annihilator computation retains the
different in every ramified factor.
\end{proof}

The density need not be locally constant on the component-vanishing locus.
For the nonvanishing outputs that occur in the global field, the shell
inverses nevertheless give a finite pointwise expansion.

\begin{lemma}[Pointwise inverse expansion]\label{lem:fourier-inverse-expansion}
If $z=(z_\nu)\in E$ has $z_\nu\ne0$ for every $\nu$, then
\begin{equation}\label{eq:fourier-inverse-expansion}
 f_p^*(z)=\sum_{\mathfrak h_p}K_{p,\mathfrak h_p}(z),
\end{equation}
where the sum is over integral $R_E$-ideals and only finitely many summands
are nonzero.
\end{lemma}

\begin{proof}
Let $B_n=p^{-n}R_E$.  Fubini and additive orthogonality give
\[
 \int_{B_n}\gamma_p(A)\psi_p(-\Tr_{E/\Q_p}(Az))\,dA
       =\vol(B_n)\,\mu_p(z+p^n\mathfrak d_E^{-1}).
\]
The self-dual normalizations give
\[
 B_n^\perp=p^n\mathfrak d_E^{-1},\qquad
 \vol(B_n)\vol(B_n^\perp)=1.
\]
Thus the right side is the average density on that neighborhood.  By
Lemma~\ref{lem:fourier-density} it equals $f_p^*(z)$ for sufficiently large
$n$.  The ball $B_n$ is the disjoint union of the finitely many shells with
$h_\nu\le ne_\nu$.

There can in fact be only finitely many nonzero shell inverses at this $z$.
The injection
$\mathfrak h_p/(\mathfrak h_p\cap\mathcal C_L)\hookrightarrow R_E/\mathcal C_L$
gives $g_p\le v_p([R_E:\mathcal C_L])$.  In
Equation~\eqref{eq:fourier-inverse-support}, a nonzero value requires
\[
 h_\nu\le v_\nu(z_\nu)+d_\nu+4e_\nu(g_p+1).
\]
In Equation~\eqref{eq:fourier-boundary-support}, it requires
$h_\nu\le v_\nu(z_\nu)+d_\nu+1$ whenever $h_\nu>0$.
All right sides are bounded for this $z$, because its component valuations
are finite.  The shell sum therefore stabilizes, and the preceding ball
inversion proves Equation~\eqref{eq:fourier-inverse-expansion}.
\end{proof}

The exact scalar coefficients in Lemma~\ref{lem:fourier-scalar-density}
supply the trace marginal in the next section.  The nonscalar estimate uses
the pointwise decay in Equation~\eqref{eq:fourier-shell-decay}, the inverse
supports in Equations~\eqref{eq:fourier-inverse-support} and
\eqref{eq:fourier-boundary-support}, and the boundary bound in
Equation~\eqref{eq:fourier-boundary-decay}.  In the case $g_p=0$ and $p>5$,
the last upper bound contains a factor $q_\nu^{-1/2}$ for each actual boundary
component where the support permits one fewer power of $\mathfrak q_\nu$.
Once the real boxes are specified, Section~\ref{sec:additive} combines the
supports with the nonzero field norm to restrict the denominator profiles in
Poisson summation.
